\section{Kinetic Voronoi Mesh}

Gemini recommends \cite{basch1999data},\\
(1999) ``Data Structures for Mobile Data''
\begin{outline}
\1 Intro / Abstract
\2 ``A kinetic data structure KDS maintains an attribute of interest in a system of
geometric objects undergoing continuous motion.''
\2 Based on partial, updating flightplans
\2 Datastructures must maintain ``certificates'' to be certified correct, like the incircle test for delaunay triangulations
\2 Foundation Paper
\end{outline}

(2004) ``An empirical comparison of techniques for updating Delaunay
triangulations'' \cite{guibas2004empirical}
\begin{outline}
\1 Abstract / Introduction
\2 Look at methods of updating delaunay meshes, what we want really.
\2 Has a great summary on kinetic data Structures 
\2 also lots of useful leads on other places this tech is used.
\2 Focused on molecular dynamics or somesuch?
\end{outline}

(2012) ``Delaynay Mesh Generation'' book
\begin{outline} \cite{cheng2013delaunay}
\1 Seems like an important resource, Note that Shewchuk made the geometric predicates we're interested in.
\1 Gemini thinks this will have info on accurate determinent predicates we need for certificates.
\1 May want to aquire this book.
\end{outline}

(2020) ``Kinetic Shape Reconstruction'' \cite{Bauchet2020kin}
\begin{outline}
\1 Gemini thinks said that CGAL had a kinetic data structure implemented already that could represent some architectural insights. 
\2 That was implemented for this paper.
\2 Solves a different problem though, building a polygon partition to mesh point clouds?
\end{outline}


